\section{Eliminating a period at controlled height}\label{sec:compression}
Instead of approximating $\zet$, we can cancel its coefficient exactly.
For a symmetric pencil $R+XB+ZC$, restriction by a rational matrix $U$
eliminates $Z$ when $U^TCU=0$. Such a subspace is called totally
isotropic for $C$. If $U$ has independent columns, the actual compressed
Gram matrix remains positive definite. The issue is constructing enough
directions without paying an excessive height cost.

\subsection{A larger cofactor space and the local cancellation}
For this construction use the full basis
\[
 E_l(u)=\frac{D(u)}{u+l^2},\quad 1\leq l\leq2K,
 \qquad \mathcal G_{lk}(X,Z)=\nu(E_lE_k/D).
\]
It is important to include polynomial parts in this matrix. For $l\ne k$,
$E_lE_k/D$ is an integer polynomial. On the diagonal, monic division
gives
\begin{equation}\label{eq:quotient}
 E_l^2/D=Q_l(u)+D'(-l^2)/(u+l^2),\qquad Q_l\in\Z[u],
\end{equation}
where
\begin{equation}\label{eq:Dprime}
 D'(-l^2)=\frac{(-1)^{l-1}(2K-l)!(2K+l)!}{2l^2}.
\end{equation}
Integration by parts in \eqref{eq:functional} gives
\begin{equation}\label{eq:moments}
 \nu(u^e)=\tfrac14\binom{2e+4}{3}T_{2e+1},
 \qquad \tan z=\sum_{e\geq0}T_{2e+1}\frac{z^{2e+1}}{(2e+1)!}.
\end{equation}
For example the first moments are $1,10,224$. One derivation expands
$g$ in exponentials after transferring three derivatives to $y^{2e+4}$;
the resulting odd-power sum is the classical tangent-number formula.
The tangent numbers are integers, also directly from
$(\tan z)'=1+\tan^2z$. Thus all polynomial moments have denominator
dividing four. Equations \eqref{eq:formal-resolvents} and
\eqref{eq:quotient} show that $32\lcm(1,\ldots,2K)^4$ clears every
coefficient of $\mathcal G$.

Both period-coefficient matrices of $\mathcal G$ are diagonal. The
$Z$ coefficient vanishes exactly on the $K$ odd coordinates, and on
the even coordinates it equals
\begin{equation}\label{eq:fullC}
 \mathcal C_{2m,2m}
 =\frac7{256}(-1)^m m(2K-2m)!(2K+2m)!.
\end{equation}
To choose an isotropic direction, first prescribe polynomial values
rather than cofactor coefficients. For $\deg P<2K$, put
$\xi_m=m^2P(-4m^2)$. Partial fractions and \eqref{eq:Dprime} give
\[
 [Z]\nu(P^2/D)=\frac7{4(4K)!}
 \sum_{m=1}^K(-1)^m m\binom{4K}{2K-2m}\xi_m^2.
\]
We will construct a zero of this quadratic form on each five-node block,
then lift a proportional evaluation vector to an integral polynomial.
Alternating signs permit cancellation, but real indefiniteness alone
does not ensure a rational zero. We use the classical Meyer theorem
in dimension five, followed by Cassels' quantitative small-zero bound
\cite[Introduction and Theorem~1.1]{Fukshansky}.

Take five consecutive indices $a\leq m\leq b=a+4\leq K$ and set
\begin{equation}\label{eq:Blocal}
 B_m^{\rm loc}=
 \frac{(2K-2a)!}{(2K-2m)!}
 \frac{(2K+2b)!}{(2K+2m)!}.
\end{equation}
These are products of exactly eight positive integers at most $4K$.
The integral diagonal form
\[
 \sum_{m=a}^b(-1)^m mB_m^{\rm loc}z_m^2
\]
is nonsingular and indefinite. Meyer supplies a rational zero, and
Cassels supplies a nonzero integer zero satisfying
\begin{equation}\label{eq:Zheight}
 \max_m|z_m|\leq[15K(4K)^8]^2=:Z_K.
\end{equation}
The useful step is removal of a common factor before applying this
bound. In evaluation coordinates the coefficients are proportional to
$(-1)^m m\binom{4K}{2K-2m}$. Within the block their common positive
rational factor is
$(4K)!/((2K-2a)!(2K+2b)!)$. Dividing by it leaves precisely the
polynomial-size coefficients in \eqref{eq:Blocal}. An estimate using
the binomial coefficients themselves would miss this saving.

For illustration, at $K=5$ the primitive equation is
\[
 -25194z_1^2+15504z_2^2-2907z_3^2+152z_4^2-z_5^2=0.
\]
The vector $(-42,57,136,561,0)$ is an exact zero. The general
existence and height statement comes from Meyer and Cassels, not from
extrapolating this example.

\subsection{An integral lift and its determinant cost}
Lift a block zero to the polynomial
\begin{equation}\label{eq:full-lift}
 P_z(u)=\sum_{m=a}^b z_mB_m^{\rm loc}E_{2m}(u).
\end{equation}
It vanishes at every odd node and at every even node outside its block.
At an even node in the block, \eqref{eq:Dprime} gives
\[
 m^2P_z(-4m^2)
 =-\tfrac18(2K-2a)!(2K+2b)!z_m.
\]
Equivalently, substituting \eqref{eq:fullC} gives its $Z$ coefficient as
$\tfrac7{256}(2K-2a)!(2K+2b)!\sum_m(-1)^m mB_m^{\rm loc}z_m^2=0$.
The cofactor coefficients are integers bounded by
\begin{equation}\label{eq:Wfull}
 W_K^{\rm full}=225K^2(4K)^{24}.
\end{equation}

\begin{theorem}\label{thm:compression}
For every $K\geq5$, there is an integral cofactor-coordinate matrix
$\mathcal U$ of rank $h=K+\lfloor K/5\rfloor$ such that
$\mathcal U^T\mathcal C\mathcal U=0$ and
\[
 \det(\mathcal U^T\mathcal U)
 \leq\bigl(5(W_K^{\rm full})^2\bigr)^{\lfloor K/5\rfloor}.
\]
Its compressed matrix is positive definite at $\bet$ and has only
rational and $\bet$ entries. Its determinant polynomial has degree
exactly $K$.
\end{theorem}
\begin{proof}
Adjoin the $K$ odd coordinate vectors to one lift
\eqref{eq:full-lift} from each of the disjoint blocks
$[5i+1,5i+5]$. Their supports are disjoint. They are independent,
mutually $\mathcal C$-orthogonal, and isotropic. The volume bound is
the product of their squared lengths. Positivity follows by restriction
of the actual Gram matrix. The beta-coefficient matrix is nonsingular
on the odd coordinates and zero on the added coordinates, so its rank
is exactly $K$. The degree assertion also follows from the block
factorization below.
\end{proof}

This is a height statement in the declared cofactor basis. The large
monomial coefficients of the cofactors have not been bounded by a
power of $K$. To see the precise transfer, let $I$ range over the
possibly nonzero row minors of $\mathcal U$: all odd coordinates and
one even coordinate from each block. Cauchy--Binet and Gram
Cauchy--Schwarz give
\begin{equation}\label{eq:transfer}
 0<\det(\mathcal U^T\mathcal G\mathcal U)(\bet)
 \leq \max_I\det\mathcal G_{I,I}(\bet,\zet)
              \bigl(25(W_K^{\rm full})^2\bigr)^{\lfloor K/5\rfloor}.
\end{equation}
Indeed the sum of the absolute row minors is the product of the
five-entry $\ell^1$ norms. The logarithm of its square is $O(K\log K)$.
If a multiplier clears every coefficient of the same mixed minors
$\det\mathcal G_{I,J}(X,Z)$, it clears the compressed polynomial too.
Thus the projection adds no leading $K^2$ height cost to a comparison
made in this same basis. It does not by itself supply that comparison.

For every fixed odd block length $L\geq5$, the extension in
Appendix~\ref{app:fixedL} retains $(L-3)/2$ independent isotropic
directions per block. The dimension becomes
$K+(L-3)\lfloor K/L\rfloor/2$, still with $O_L(K\log K)$
logarithmic height cost. Fixing $L$ before $K\to\infty$ is essential.

\subsection{The rational block and what survives it}
Order the odd coordinates first. The compressed matrix has the form
\begin{equation}\label{eq:compression-schur}
 \begin{pmatrix}L_0+XB_0&T_0\\T_0^T&E_0\end{pmatrix},
 \qquad
 \det=\det E_0\,
       \det(L_0+XB_0-T_0E_0^{-1}T_0^T).
\end{equation}
Here $T_0,E_0$ are rational, $E_0>0$, and $B_0$ is nonsingular.
The primitive polynomial therefore has degree $K$, not the full
matrix dimension. The rational factor $\det E_0$ disappears, whereas
the rational correction inside the Schur complement remains.
Consequently this reduction does not show that all benefit of the
integer poles disappears.

This also explains a useful invariance. Adding an element of the
full $Z$-radical to an extra lift is a determinant-one basis shear;
it leaves the determinant polynomial exactly unchanged. More generally,
rational changes of basis multiply it by a rational square and leave
its primitive polynomial unchanged. Denominators introduced by a
particular interpolation of the same subspace are not automatically
an intrinsic determinant cost.

At the two checked indices $K=5,10$, the logarithms of the primitive
values are respectively $101.17\ldots$ and $475.49\ldots$, compared
with $94.29\ldots$ and $444.32\ldots$ for the radical alone. These
finite comparisons motivate the lower-degree variant below; they do
not establish an asymptotic failure theorem.
